\chapter{Introdução}\label{chp:intro}

O setor elétrico brasileiro possui uma matriz predominantemente renovável, com destaque para as usinas hidrelétricas, que historicamente respondem pela maior parte da geração de energia elétrica no país. No entanto, períodos de estiagem prolongada, aliados ao crescimento da demanda e à necessidade de segurança energética, exigem o acionamento frequente de usinas termelétricas. Embora essas usinas sejam fundamentais para garantir o suprimento contínuo de energia, sua operação está associada à emissão de gases de efeito estufa, especialmente o \ac{CO$_2$}. As mudanças climáticas, impulsionadas pelas emissões de diversos setores econômicos, têm imposto desafios crescentes aos planejadores energéticos. No Brasil, o \ac{SIN} requer ferramentas que permitam equilibrar a segurança do suprimento elétrico com as metas de descarbonização assumidas pelo país em acordos internacionais. Nesse contexto, a estimação precisa das emissões de \ac{CO$_2$} das termelétricas em escala horária torna-se um insumo relevante para o despacho econômico-ambiental e para o planejamento da operação.

A literatura internacional apresenta modelos consolidados para a estimação de emissões horárias de usinas termelétricas, geralmente baseados em funções quadráticas ou exponenciais calibradas a partir de medições contínuas ou de dados de projeto de cada usina. No entanto, o cenário brasileiro apresenta uma particularidade: os dados de emissões de \ac{CO$_2$} das termelétricas são inventariados anualmente — e não horariamente — por instituições como o \ac{IEMA} e, anteriormente, pela Eletrobras. Essa disponibilidade exclusivamente anual de dados impõe um problema de natureza inversa: é necessário estimar o comportamento horário das emissões a partir de agregados anuais. Além disso, as termelétricas brasileiras operam com diferentes tipos de combustível (gás natural, óleo diesel, carvão mineral, biomassa), diferentes ciclos de operação (ciclo simples, ciclo combinado) e operam em patamares discretos de carga, com transições entre níveis de geração que afetam a eficiência e, consequentemente, as emissões. A literatura nacional ainda carece de metodologias sistemáticas para a discretização temporal das emissões anuais em estimativas horárias que considerem essas particularidades operacionais e que sejam generalizáveis para diferentes usinas.

Diante desse cenário, este trabalho propõe uma metodologia para estimar as emissões horárias de \ac{CO$_2$e} das usinas termelétricas brasileiras, combinando técnicas de otimização não-linear e redes neurais artificiais. A metodologia estrutura-se em duas etapas principais. Na primeira etapa, calibram-se os coeficientes de uma equação paramétrica de emissão (composta por termos quadrático e exponencial) por meio de um problema de minimização resolvido em \ac{AMPL} com o solver \ac{IPOPT}, utilizando como referência os totais anuais do \ac{IEMA}. Dois cenários são comparados: calibração padrão e calibração com regularização L2. Na segunda etapa, as emissões sintéticas horárias geradas tornam-se alvos para o treinamento de uma rede neural \ac{MLP}, que aprende diretamente a relação entre variáveis operacionais (geração, categoria de operação, fase de transição), características estáticas da usina (combustível, potência instalada, eficiência) e as emissões horárias. Mais especificamente, o trabalho desenvolve um algoritmo de classificação de categorias de geração e identificação de transições operacionais com base em limiares adaptativos, calibra os coeficientes da equação de emissão para 39 usinas termelétricas brasileiras no período de 2020 a 2024, treina e valida uma rede neural \ac{MLP} para estimar emissões horárias, e compara o desempenho dos cenários de calibração (padrão vs. L2) em termos de precisão e estabilidade.

Como contribuições principais, este trabalho oferece uma metodologia sistemática para discretização de emissões anuais em estimativas horárias aplicável a termelétricas com diferentes combustíveis e ciclos de operação; uma base de dados integrada combinando informações da \ac{ONS} e do \ac{IEMA} para 39 usinas em um período de 5 anos, com enriquecimento de \textit{features} temporais (categorias, transições, fases); uma comparação quantitativa entre calibração padrão e com regularização L2, fornecendo subsídios para a escolha da estratégia mais adequada em problemas de calibração inversa no setor elétrico; e um modelo de rede neural treinado para estimação horária de emissões, com análise de desempenho por categoria operacional e tipo de combustível. Essas contribuições podem apoiar o planejamento energético e ambiental do \ac{SIN}, fornecendo insumos horários para modelos de despacho econômico-ambiental e para a avaliação de políticas de descarbonização.